Wigner 小群已经说明：有质量粒子的内部自由度表现为自旋，而无质量粒子的物理内部自由度表现为螺旋度。散射振幅要使用这些量子数，还需要把它们写成能够直接与 Dirac 链或偏振张量缩并的投影算符。

对有质量电子，螺旋度由协变自旋四矢量沿 $\mathbf p$ 平行或反平行的两种取向给出。对真实光子，它对应二维横向偏振平面上的圆偏振本征基。改变外态偏振只改变 LSZ 外腿的投影，不改变 QED 相互作用顶角本身。


\begin{equation}
 \boxed{
 u_h(p)\bar u_h(p)
 =\frac12(\slashed p+\me c)
 \left(1+\gamma^5\slashed s_h\right).}
 \label{eq:electron-helicity-projector-dp9}
\end{equation}

对极端相对论过程，有时还需要区分螺旋度与手征性。手征投影算符已在 \eqnrefs{eq:chiral-projectors} 定义为
\begin{equation*}
 P_{R,L}=\frac12(1\pm\gamma^5).
\end{equation*}
对 $\me\neq0$，$P_{R,L}$ 不等于螺旋度投影算符；只有在严格无质量 Dirac 方程或受控 $\me c^2/E\to0$ 极限中，两者才沿 positive-能量粒子分支对齐。

\begin{equation}
 \boxed{
 \mathsf P_\lambda^{\mu\nu}(k;n)
 \equiv
 \varepsilon_\lambda^\mu\varepsilon_\lambda^{\nu *}
 =\frac12\left[
 \mathsf P_\perp^{\mu\nu}(k;n)
 -\ii\lambda\mathsf C^{\mu\nu}(k;n)
 \right].}
 \label{eq:photon-helicity-projector-covariant-dp9}
\end{equation}

若振幅写成 $\mathcal M_\lambda=\varepsilon_{\lambda\mu}\mathcal M^\mu$ 且 Ward 条件 $k_\mu\mathcal M^\mu=0$ 已满足，则固定螺旋度的概率核为
\begin{equation}
 |\mathcal M_\lambda|^2
 =\mathcal M_\mu
 \mathsf P_\lambda^{\mu\nu}
 \mathcal M_\nu^*.
 \label{eq:helicity-resolved-photon-amplitude-dp9}
\end{equation}
这里 $n^\mu$ 的改变只改变极化代表元的规范部分；对守恒的振幅，物理结果与 $n^\mu$ 无关。若这一独立性在计算中失败，应先检查 Ward 恒等式，而不能把参考矢量依赖当成可观测偏振效应。

这些投影式已经把常见极化制备统一到同一套密度矩阵语言中：费米子由 \eqnrefs{eq:polarized-spinor-projector-r9,eq:electron-helicity-projector-dp9} 描述，光子由 \eqnrefs{eq:photon-helicity-projector-covariant-dp9} 描述；自旋平均、固定自旋、圆极化或螺旋度不对称只对应不同的实验制备、求和与取差方式，并不需要额外的经验规则。

